% !PRINTED 87
\weberpage{87}
Demnach hat die Congruenz (17) zur Folge:
\begin{gather}
e_{1,h}i_{1,h}p^{\pi-\pi_1}x_1 + e_{2,h}i_{2,h}p^{\pi-\pi_2}x_2 + \cdots e_{\varrho,h}i_{\varrho,h}p^{\pi-\pi_\varrho}x_\varrho \notag\\
\equiv 0\ (\mathrm{mod}\ p^\pi) \qquad h = 1, 2, \dots \mu, \tag{19}
\end{gather}
und die übrigen in $J$ aufgehenden Primzahlen ergeben jede ein
eben solches System von Congruenzen, die in ihrer Gesammtheit
vollständig gleichbedeutend mit der Congruenz (17) sind. Diese
Congruenzen (19) dürfen nicht anders befriedigt werden können,
als dadurch, dass die Zahlen
\[
p^{\pi-\pi_1}x_1,\ p^{\pi-\pi_2}x_2,\ \dots\ p^{\pi-\pi_\varrho}x_\varrho
\]
durch $p^\pi$ theilbar sind.

    Setzen wir also zur Abkürzung:
\[
p^{\pi-\pi_1}x_1 = y_1,\ p^{\pi-\pi_2}x_2 = y_2,\ \dots\ p^{\pi-\pi_\varrho}x_\varrho = y_\varrho,
\]
so erhält (19) die Form:
\weq{20}{e_{1,h}i_{1,h}y_1 + e_{2,h}i_{2,h}y_2 + \cdots e_{\varrho,h}i_{\varrho,h}y_\varrho \equiv 0\ (\mathrm{mod}\ p^\pi),}
und diese Congruenzen sollen
\[
y_1 \equiv 0,\ y_2 \equiv 0,\ \dots\ y_\varrho \equiv 0\ (\mathrm{mod}\ p^\pi)
\]
zur nothwendigen Folge haben.

    Das System von linearen Congruenzen (20) kann nun mit
denselben Mitteln discutirt werden, wie ein System linearer Gleichungen, nämlich durch die Determinanten (vgl. den Abschnitt II
des ersten Bandes).

    Wir wollen zur besseren Uebersicht die Congruenzen (20)
für den Augenblick so schreiben:
\begin{gather}
a_{1,1}y_1 + a_{2,1}y_2 + \cdots + a_{\varrho,1}y_\varrho \equiv 0 \notag\\
a_{1,2}y_1 + a_{2,2}y_2 + \cdots + a_{\varrho,2}y_\varrho \equiv 0\ (\mathrm{mod}\ p^\pi). \tag{21}\\
\dots\dots\dots\dots\dots\dots\dots\dots\dots\dots \notag\\
a_{1,\mu}y_1 + a_{2,\mu}y_2 + \cdots + a_{\varrho,\mu}y_\varrho \equiv 0 \notag
\end{gather}
Wenn die Coëfficienten $a_{k,h}$ alle durch $p$ theilbar sind, dann
sind diese Congruenzen schon befriedigt, wenn die $y$ alle nur
durch $p^{\pi-1}$ theilbar sind, und dieser Fall darf also nicht eintreten.

    Wir wollen nun annehmen, es lasse sich aus der Matrix
\weq{22}{
\begin{matrix}
a_{1,1}, & a_{2,1}, & \dots & a_{\varrho,1} \\
a_{1,2}, & a_{2,2}, & \dots & a_{\varrho,2} \\
\dots & \dots & \dots & \dots \\
a_{1,\mu}, & a_{2,\mu}, & \dots & a_{\varrho,\mu}
\end{matrix}
}
%CONT>>
